\documentclass[11pt,a4paper]{article}

\usepackage[margin=1in]{geometry}
\usepackage{amsmath,amssymb,amsthm}
\usepackage{booktabs}
\usepackage{multirow}
\usepackage{microtype}
\usepackage{graphicx}
\emergencystretch=4em
\usepackage{hyperref}
\hypersetup{colorlinks=true,linkcolor=black,urlcolor=blue,citecolor=black}

\theoremstyle{plain}
\newtheorem{proposition}{Proposition}
\newtheorem{lemma}[proposition]{Lemma}
\newtheorem{corollary}[proposition]{Corollary}
\theoremstyle{definition}
\newtheorem{remark}[proposition]{Remark}

\newcommand{\R}{\mathbb{R}}
\newcommand{\E}{\mathcal{I}}
\newcommand{\tr}{\operatorname{tr}}
\newcommand{\diag}{\operatorname{diag}}
\newcommand{\spec}{\operatorname{spec}}
\newcommand{\supp}{\operatorname{supp}}
\newcommand{\vecop}{\operatorname{vec}}
\newcommand{\cond}{\kappa}

\title{Closed-Form Observed Information for Coordinate-Based Meta-Regression}
\author{Derivations, equivalence proofs, and the structure of the covariance step}
\date{25 August 2026}

\begin{document}
\maketitle

\begin{abstract}
NiMARE's coordinate-based meta-regression (CBMR) computes regression standard errors by
differentiating the log-likelihood twice with \texttt{torch.func.hessian}. Forward-over-reverse
mode pushes every tangent through the likelihood at once, so the intermediate has shape
$(\text{parameters} \times \text{patterns} \times \text{voxels})$; at $7{,}329$ coefficients,
$21$ patterns and $8{,}444$ voxels that is $9.7$~GB, which is an out-of-memory failure rather
than a slow fit.

We derive the observed information in closed form for all three of NiMARE's observation
distributions, prove each derivation equivalent to the autodiff Hessian, and characterise the
sparsity the design induces in it. A single lemma covers all three: the second differential of
the log-likelihood with respect to the log-intensity is either diagonal in the voxels (Poisson,
Negative Binomial) or diagonal-plus-rank-one (Clustered Negative Binomial), and the chain rule
through the log-bilinear predictor turns each into $n_{\text{patterns}}$ symmetric rank-updates
of width $n_{\text{bases}}$. No parameter-sized intermediate is ever formed. We then separate
two questions that are easily conflated: which \emph{likelihood} admits a closed-form Hessian
(all three), and which \emph{designs} make that Hessian block diagonal (a strictly narrower
class), and we give the generalisations that recover the fast covariance step outside it.
\end{abstract}

\tableofcontents

\section{Setting and notation}

\subsection{The pattern representation}

CBMR models the foci of $n$ experiments over $V$ voxels as a point process whose log-intensity
is bilinear in a spatial spline basis and an experiment-level design. Write

\begin{itemize}
  \item $B \in \R^{V \times K}$ for the spline basis, $b_v^\top$ its $v$-th row;
  \item $L \in \R^{P \times C}$ for the \emph{loadings}: the distinct rows of the experiment-level
        blocks of every spatial term, stacked. $P$ is the number of distinct spatial loadings
        (``patterns''), $C$ the number of spatial columns;
  \item $\pi : \{1,\dots,n\} \to \{1,\dots,P\}$ for the pattern of each experiment, and
        $\E_p = \pi^{-1}(p)$;
  \item $\beta \in \R^{C \times K}$ for the spatial coefficients;
  \item $G \in \R^{n \times Q}$ for the non-spatial (``global'') design block and
        $\gamma \in \R^{Q}$ for its coefficients.
\end{itemize}

The log-intensity of pattern $p$ at voxel $v$, and the scalar moderator effect of experiment
$i$, are
\begin{equation}\label{eq:predictor}
  S_{pv} \;=\; \sum_{c=1}^{C} L_{pc} \sum_{k=1}^{K} \beta_{ck} B_{vk},
  \qquad
  m_i \;=\; \sum_{q=1}^{Q} G_{iq}\gamma_q ,
\end{equation}
so that $S = (L\beta)B^\top$ and $m = G\gamma$. We write $u_{pv} = e^{S_{pv}}$ and
$e_i = e^{m_i}$ throughout. Foci counts are $y_{iv}$, with pattern marginals and experiment
totals
\begin{equation}
  Y_{pv} \;=\; \sum_{i \in \E_p} y_{iv},
  \qquad
  y_{i\cdot} \;=\; \sum_{v=1}^{V} y_{iv}.
\end{equation}

The parameter vector is $\phi = (\vecop(\beta), \gamma) \in \R^{CK + Q}$, with $\beta$ flattened
row-major so that $\beta_{ck}$ occupies index $cK + k$. Overdispersion parameters $\theta_p$,
where the distribution has them, are held fixed at their fitted values --- this is what
\texttt{CBMRModel.information\_matrix} does, and it is what CBMR has always reported regression
standard errors from. Every derivation below is therefore a differentiation in $\phi$ only,
with $\theta$ a constant.

\subsection{The two facts the whole paper rests on}

Both parts of \eqref{eq:predictor} are \emph{linear} in the parameters. Hence
\begin{equation}\label{eq:linear}
  \frac{\partial S_{pv}}{\partial \beta_{ck}} = L_{pc} B_{vk},
  \qquad
  \frac{\partial^2 S_{pv}}{\partial \beta_{ck}\,\partial \beta_{c'k'}} = 0,
  \qquad
  \frac{\partial m_i}{\partial \gamma_q} = G_{iq},
  \qquad
  \frac{\partial^2 m_i}{\partial \gamma_q\,\partial \gamma_{q'}} = 0 .
\end{equation}
The second derivatives vanish identically. Everything that follows is the chain rule applied to
this fact, so the only distribution-specific work is the second differential of the
log-likelihood with respect to $(S, \gamma)$.

\section{A master lemma}

Every CBMR log-likelihood in NiMARE decomposes over patterns,
\begin{equation}\label{eq:decomposition}
  -\log \mathcal{L}(\phi) \;=\; \sum_{p=1}^{P} f_p\bigl(S_{p\cdot},\, \gamma\bigr) ,
\end{equation}
because experiments sharing a pattern share a log-intensity map, and no term couples two
patterns. This is the representation \texttt{predictor.py} is built on.

\begin{lemma}[Chain rule through a log-bilinear predictor]\label{lem:master}
Let $F(\phi) = \sum_p f_p(S_{p\cdot}, \gamma)$ with $S$ as in \eqref{eq:predictor} and each $f_p$
twice differentiable. Write
\[
  \bigl(\Sigma_p\bigr)_{vv'} = \frac{\partial^2 f_p}{\partial S_{pv} \partial S_{pv'}},
  \qquad
  \bigl(\tau_p\bigr)_{vq} = \frac{\partial^2 f_p}{\partial S_{pv} \partial \gamma_q},
  \qquad
  \bigl(\Xi_p\bigr)_{qq'} = \frac{\partial^2 f_p}{\partial \gamma_q \partial \gamma_{q'}} .
\]
Then the Hessian of $F$ has blocks
\begin{align}
  \frac{\partial^2 F}{\partial \beta_{ck}\,\partial \beta_{c'k'}}
    &= \sum_{p=1}^{P} L_{pc} L_{pc'} \, \bigl(B^\top \Sigma_p B\bigr)_{kk'},
      \label{eq:master-bb}\\
  \frac{\partial^2 F}{\partial \beta_{ck}\,\partial \gamma_q}
    &= \sum_{p=1}^{P} L_{pc} \, \bigl(B^\top \tau_p\bigr)_{kq},
      \label{eq:master-bg}\\
  \frac{\partial^2 F}{\partial \gamma_q\,\partial \gamma_{q'}}
    &= \sum_{p=1}^{P} \bigl(\Xi_p\bigr)_{qq'} .
      \label{eq:master-gg}
\end{align}
\end{lemma}

\begin{proof}
By \eqref{eq:linear}, $\partial S_{pv}/\partial\beta_{ck} = L_{pc}B_{vk}$ and the second
derivative of $S$ vanishes, so the chain rule for a composition with an affine map contributes
no curvature term:
\[
  \frac{\partial^2 f_p}{\partial \beta_{ck} \partial \beta_{c'k'}}
  = \sum_{v,v'} \frac{\partial^2 f_p}{\partial S_{pv}\partial S_{pv'}}
      \frac{\partial S_{pv}}{\partial\beta_{ck}} \frac{\partial S_{pv'}}{\partial\beta_{c'k'}}
  = L_{pc}L_{pc'} \sum_{v,v'} (\Sigma_p)_{vv'} B_{vk}B_{v'k'} ,
\]
which is \eqref{eq:master-bb}. The same argument with one factor replaced by
$\partial \gamma_q/\partial\gamma_q = 1$ gives \eqref{eq:master-bg}, and with both replaced,
\eqref{eq:master-gg}. Summing over $p$ completes the proof.
\end{proof}

Lemma~\ref{lem:master} reduces each distribution to a single question: \emph{what is $\Sigma_p$?}
The answer determines the cost. If $\Sigma_p = \diag(w_p)$ then
$B^\top \Sigma_p B = B^\top \diag(w_p) B$ costs $O(VK^2)$ and allocates one $V \times K$
temporary. If $\Sigma_p$ is diagonal plus rank one, one extra outer product of a
$K$-vector is added. In neither case is anything of size $CK$ formed.

\begin{remark}
This is exactly what forward-over-reverse autodiff fails to exploit. \texttt{jacfwd} carries one
tangent per parameter through the evaluation, materialising an array of shape $(CK, P, V)$,
because it cannot know that the map $\beta \mapsto S$ is affine with a rank structure that
collapses the parameter axis.
\end{remark}

\section{Poisson}

\subsection{The likelihood}

Dropping the parameter-free $-\sum \log y!$, NiMARE's Poisson log-likelihood on marginals is
\begin{equation}\label{eq:poisson-ll}
  \log\mathcal{L} \;=\; \sum_{p}\sum_{v} Y_{pv} S_{pv}
  \;+\; \sum_{i} y_{i\cdot}\, m_i
  \;-\; \sum_{p} \Bigl(\sum_v u_{pv}\Bigr) T_p,
  \qquad T_p = \sum_{i \in \E_p} e_i .
\end{equation}

\begin{proposition}[Poisson observed information]\label{prop:poisson}
With $U_{pq} = \sum_{i\in\E_p} e_i G_{iq}$, $a_{pk} = \sum_v u_{pv}B_{vk}$ and
$E_p = \sum_v u_{pv}$, the Hessian of $-\log\mathcal{L}$ is
\begin{align}
  H_{(ck),(c'k')} &= \sum_p L_{pc}L_{pc'} \bigl(B^\top \diag(T_p u_{p\cdot}) B\bigr)_{kk'},\\
  H_{(ck),q} &= \sum_p L_{pc}\, a_{pk} U_{pq},\\
  H_{q,q'} &= \sum_i e_i E_{\pi(i)} G_{iq}G_{iq'} .
\end{align}
In particular $H$ does not depend on the foci.
\end{proposition}

\begin{proof}
Write \eqref{eq:poisson-ll} in the form \eqref{eq:decomposition} with
$f_p = -\sum_v Y_{pv}S_{pv} - \sum_{i\in\E_p} y_{i\cdot}m_i + E_p(S_{p\cdot})\,T_p(\gamma)$.
The first two terms are linear in $\phi$ and so contribute nothing to any second derivative;
this is the whole of the data dependence, which proves the final claim. For the remainder,
$\partial E_p/\partial S_{pv} = u_{pv}$ and $\partial^2 E_p/\partial S_{pv}\partial S_{pv'} =
u_{pv}\delta_{vv'}$, while $\partial T_p/\partial\gamma_q = U_{pq}$ and
$\partial^2 T_p / \partial\gamma_q\partial\gamma_{q'} = \sum_{i\in\E_p} e_i G_{iq}G_{iq'}$.
Hence $\Sigma_p = T_p \diag(u_{p\cdot})$, $(\tau_p)_{vq} = u_{pv}U_{pq}$ and
$(\Xi_p)_{qq'} = E_p \sum_{i \in \E_p} e_iG_{iq}G_{iq'}$. Substituting into
Lemma~\ref{lem:master} and collecting the $\gamma\gamma$ sum over experiments gives the three
displays.
\end{proof}

\begin{remark}
The foci-independence is the log link's canonical-parameter property: for a GLM with canonical
link the observed and expected information coincide. It is a useful test of an implementation,
since a Poisson Hessian that responds to the data is wrong; the test suite checks this.
\end{remark}

\section{Negative Binomial}

\subsection{The likelihood in $(\eta_p, \kappa_p)$}

NiMARE's negative binomial is a moment-matched approximation to the sum, over the experiments of
a pattern, of negative binomial variables. With overdispersion $\theta_p$ and
\begin{equation}
  s_{1} = \sum_{i\in\E_p} e_i, \qquad s_{2} = \sum_{i\in\E_p} e_i^2 ,
\end{equation}
it uses shape $r = s_1^2/(\theta_p s_2)$ and success probability
$\pi_v = \bigl(1 + s_1/(\theta_p u_{pv} s_2)\bigr)^{-1}$. That is, $Y_{pv}$ is negative binomial
with mean $\mu_{pv} = s_1 u_{pv}$ and dispersion $k_p = \theta_p s_2/s_1^2$, so that
$\operatorname{Var} Y_{pv} = \mu_{pv} + k_p\mu_{pv}^2$ and $r = 1/k_p$. We work in
\begin{equation}\label{eq:nb-mk}
  \eta_p \;=\; \log s_1, \qquad \kappa_p \;=\; \log k_p, \qquad x_{pv} \;=\; k_p\mu_{pv},
\end{equation}
which give $\pi_v = x_{pv}/(1+x_{pv})$ and $1-\pi_v = 1/(1+x_{pv})$.

\begin{lemma}\label{lem:nb-lambda}
Up to an additive constant depending only on the data,
$f_p = -\sum_v \lambda(\eta_p, \kappa_p, S_{pv}; Y_{pv})$ with
\begin{equation}\label{eq:lambda}
  \lambda(\eta, \kappa, S; Y) \;=\; \log\Gamma(Y+e^{-\kappa}) - \log\Gamma(e^{-\kappa})
     + Y\kappa + Y(\eta+S) - (Y+e^{-\kappa})\log\bigl(1+e^{\kappa+\eta+S}\bigr).
\end{equation}
\end{lemma}

\begin{proof}
Substituting $r = e^{-\kappa}$, $\log(1-\pi_v) = -\log(1+x_{pv})$ and
$\log \pi_v = \kappa_p + \eta_p + S_{pv} - \log(1+x_{pv})$ into NiMARE's summand
$\log\Gamma(Y+r) - \log\Gamma(Y+1) - \log\Gamma(r) + r\log(1-\pi) + Y\log\pi$ gives
\eqref{eq:lambda} minus $\log\Gamma(Y_{pv}+1)$, which is free of $\phi$.
\end{proof}

The value of \eqref{eq:lambda} is that it separates the three ways the parameters enter: the
special functions see only $\kappa_p$, the voxels enter only through $\eta_p + S_{pv}$, and the
moderators only through $(\eta_p, \kappa_p)$. Every derivative in $S_{pv}$ is therefore the
corresponding derivative in $\eta_p$.

\begin{proposition}[Negative binomial observed information]\label{prop:nb}
Write
\begin{equation}\label{eq:nb-sigma}
  w_{pv} \;=\; \frac{\mu_{pv} + Y_{pv}x_{pv}}{(1+x_{pv})^{2}}, \qquad
  c_{pv} \;=\; \frac{(Y_{pv} - \mu_{pv})\,x_{pv}}{(1+x_{pv})^{2}}, \qquad
  g(x) \;=\; \frac{\log(1+x)}{x} - \frac{1}{1+x},
\end{equation}
and, for integer counts, with $n_{pj}$ the number of voxels at which $Y_{pv} > j$,
\begin{equation}\label{eq:nb-count-sums}
  G_{1p} = \sum_{j\ge 0} n_{pj}\frac{jk_p}{1+jk_p}, \qquad
  G_{2p} = \sum_{j\ge 0} n_{pj}\frac{jk_p}{(1+jk_p)^2}.
\end{equation}
Then $\Sigma_p = \diag(w_{p\cdot})$, $(\tau_p)_{vq} = w_{pv}\,\partial_q\eta_p + c_{pv}\,\partial_q\kappa_p$,
and
\begin{equation}\label{eq:nb-gg}
\begin{split}
  (\Xi_p)_{qq'} \;=\;{}& \Bigl(\sum_v w_{pv}\Bigr)\partial_q\eta_p\,\partial_{q'}\eta_p
    + \Bigl(\sum_v c_{pv}\Bigr)\bigl(\partial_q\eta_p\,\partial_{q'}\kappa_p
      + \partial_{q'}\eta_p\,\partial_q\kappa_p\bigr)
    - \Lambda_{\kappa\kappa}\,\partial_q\kappa_p\,\partial_{q'}\kappa_p\\
    &{}- \Lambda_\eta\,\partial^2_{qq'}\eta_p - \Lambda_\kappa\,\partial^2_{qq'}\kappa_p ,
\end{split}
\end{equation}
where
\begin{align}
  \Lambda_\eta &= \sum_v \frac{Y_{pv}-\mu_{pv}}{1+x_{pv}}, \notag\\
  \Lambda_\kappa &= G_{1p} + \sum_v \Bigl(\mu_{pv}\,g(x_{pv})
     - \frac{Y_{pv}x_{pv}}{1+x_{pv}}\Bigr), \label{eq:nb-second}\\
  \Lambda_{\kappa\kappa} &= G_{2p} + \sum_v \Bigl(\mu_{pv}\Bigl(\frac{x_{pv}}{(1+x_{pv})^2}
     - g(x_{pv})\Bigr) - \frac{Y_{pv}x_{pv}}{(1+x_{pv})^2}\Bigr). \notag
\end{align}
\end{proposition}

\begin{proof}
Differentiate \eqref{eq:lambda}, using $\partial\mu/\partial\eta = \mu$, $\partial x/\partial\eta =
\partial x/\partial\kappa = x$ and $e^{-\kappa}x = \mu$:
\[
  \lambda_\eta = Y - \frac{(Y+e^{-\kappa})x}{1+x} = \frac{Y-\mu}{1+x}, \qquad
  \lambda_{\eta\eta} = -\frac{\mu+Yx}{(1+x)^2}, \qquad
  \lambda_{\eta\kappa} = -\frac{(Y-\mu)x}{(1+x)^2}.
\]
Since $\lambda$ is a sum of per-voxel terms and $S_{pv}$ enters as $\eta_p$ does, $\Sigma_p$ is
diagonal with entries $-\lambda_{\eta\eta} = w_{pv}$, and $(\tau_p)_{vq}$ is
$-\lambda_{\eta\eta}\partial_q\eta_p - \lambda_{\eta\kappa}\partial_q\kappa_p$. In $\kappa$,
\[
  \lambda_\kappa = Y - e^{-\kappa}\bigl(\psi(Y+e^{-\kappa}) - \psi(e^{-\kappa})\bigr)
     + \mu\,g(x) - \frac{Yx}{1+x}.
\]
For integer $Y$, $\psi(Y+r) - \psi(r) = \sum_{j<Y} 1/(r+j)$, so the first two terms are
$\sum_{j<Y} jk/(1+jk)$; summed over voxels this is $G_{1p}$. Differentiating once more, and using
$\psi_1(Y+r) - \psi_1(r) = -\sum_{j<Y} 1/(r+j)^2$, gives $\Lambda_{\kappa\kappa}$ with $G_{2p}$.
Since $\gamma$ enters $f_p$ only through $(\eta_p, \kappa_p)$, the chain rule gives
\eqref{eq:nb-gg}, the last two terms being the curvature of the inner map.
\end{proof}

\begin{lemma}[Derivatives of $\eta_p$ and $\kappa_p$]\label{lem:nb-mk-derivs}
Let $U_q = \sum_{i\in\E_p} e_iG_{iq}$, $W_q = \sum_{i\in\E_p} e_i^2G_{iq}$,
$U^{(2)}_{qq'} = \sum_{i\in\E_p} e_iG_{iq}G_{iq'}$,
$W^{(2)}_{qq'} = \sum_{i\in\E_p} e_i^2G_{iq}G_{iq'}$, and
\[
  \ell_q = \frac{U_q}{s_1},\quad
  h_q = \frac{2W_q}{s_2},\quad
  \ell^{(2)}_{qq'} = \frac{U^{(2)}_{qq'}}{s_1} - \frac{U_qU_{q'}}{s_1^2},\quad
  h^{(2)}_{qq'} = \frac{4W^{(2)}_{qq'}}{s_2} - \frac{4W_qW_{q'}}{s_2^2} .
\]
Then
\[
  \partial_q \eta_p = \ell_q, \qquad
  \partial_q \kappa_p = h_q - 2\ell_q, \qquad
  \partial^2_{qq'} \eta_p = \ell^{(2)}_{qq'}, \qquad
  \partial^2_{qq'} \kappa_p = h^{(2)}_{qq'} - 2\ell^{(2)}_{qq'} .
\]
\end{lemma}

\begin{proof}
$\eta_p = \log s_1$ and $\kappa_p = \log\theta_p + \log s_2 - 2\log s_1$. Since
$\partial_q s_1 = U_q$ and $\partial_q s_2 = 2W_q$, we get $\partial_q \log s_1 = \ell_q$ and
$\partial_q \log s_2 = h_q$; differentiating again, $\partial^2_{qq'}\log s_1 = \ell^{(2)}_{qq'}$
and $\partial^2_{qq'} \log s_2 = h^{(2)}_{qq'}$ (the factor $4$ arises because
$\partial^2_{qq'} s_2 = 4W^{(2)}_{qq'}$).
\end{proof}

\begin{remark}[Why $(\eta, \kappa)$]\label{rem:nb-precision}
An earlier form of this derivation worked in $R_p = s_1^2/(\theta_p s_2) = 1/k_p$ and
$A_p = s_1/(\theta_p s_2)$. It is the same mathematics, but not the same arithmetic. As
$\theta_p \to 0$ both grow like $1/\theta_p$, and the moderator block subtracts terms of that
order that should cancel to order one. On a $1{,}328$-experiment voxel-based meta-regression,
one pattern of $99$ experiments fitted $\theta_p = 8.4\times10^{-13}$, so $r \approx 1.6\times10^{12}$.
For that pattern's diagonal moderator entry a $50$-digit evaluation gives $30.92$; the $(R,A)$
form gave $5.89$, and the information had a negative eigenvalue. In $(\eta,\kappa)$ every term
of Proposition~\ref{prop:nb} is of order one: $w$, $c$ and the summands of $G_{1p}, G_{2p}$
are bounded, and $g(x) \approx x/2$ is summed from its power series
$\sum_{n\ge1} (-1)^{n+1} n x^n/(n+1)$ below $x = 10^{-2}$, where the direct form cancels. The
$(\eta,\kappa)$ form agrees with the $50$-digit value to $10^{-7}$ relative. Autodiff gave
$27.69$: it differentiates the likelihood as NiMARE writes it, $\log\Gamma(Y+r) - \log\Gamma(r)$
with $r \approx 10^{12}$, which is not exact there either.

The sums \eqref{eq:nb-count-sums} need integer counts. CBMR's foci are integers, so their
pattern marginals are too; for fractional counts the implementation uses the digamma forms,
which are exact away from $k_p = 0$.
\end{remark}

\begin{corollary}[Poisson as a limit]
As $\theta_p \to 0$ with $\gamma$ fixed, $x_{pv} \to 0$ and
\[
  w_{pv} = \frac{\mu_{pv} + Y_{pv}x_{pv}}{(1+x_{pv})^2}
   \;\longrightarrow\; s_1 u_{pv} \;=\; T_p u_{pv},
\]
recovering $\Sigma_p$ of Proposition~\ref{prop:poisson}. In particular the negative binomial
weight is foci-dependent at any $\theta_p > 0$ and the dependence vanishes only in the Poisson
limit. In the same limit $c_{pv}$, $G_{1p}$, $G_{2p}$ and $\mu_{pv}g(x_{pv})$ vanish, so
\eqref{eq:nb-gg} tends to the information of a Poisson likelihood on the pattern marginals. That
differs from the $\gamma\gamma$ block of Proposition~\ref{prop:poisson}, whose count term
$\sum_i y_{i\cdot}m_i$ is taken per experiment, by $-\bigl(\sum_v Y_{pv}\bigr)\ell^{(2)}_{qq'}$.
\end{corollary}

\section{Clustered Negative Binomial}

\subsection{The likelihood}

Here the latent effect is a property of the experiment rather than of the voxel, so the spatial
coefficients enter only through the \emph{scalar} $E_p = \sum_v u_{pv}$. With precision
$\nu_p = 1/\theta_p$ and $n_p = |\E_p|$,
\begin{equation}\label{eq:cnb-ll}
\begin{split}
  \log\mathcal{L}_p ={}& n_p \nu_p \log \nu_p - n_p\log\Gamma(\nu_p)
   + \sum_{i\in\E_p}\log\Gamma(y_{i\cdot}+\nu_p)\\
   &{}- \sum_{i\in\E_p}(y_{i\cdot}+\nu_p)\log\bigl(E_p e_i + \nu_p\bigr)
   + \sum_v Y_{pv}S_{pv} + \sum_{i\in\E_p} y_{i\cdot} m_i .
\end{split}
\end{equation}

\begin{proposition}[Clustered negative binomial observed information]\label{prop:cnb}
Let $D_i = E_{\pi(i)} e_i + \nu_{\pi(i)}$, $a_{pk} = \sum_v u_{pv}B_{vk}$ and
\[
  g'_p = \sum_{i\in\E_p} \frac{(y_{i\cdot}+\nu_p)e_i}{D_i},
  \qquad
  g''_p = -\sum_{i\in\E_p} \frac{(y_{i\cdot}+\nu_p)e_i^2}{D_i^2} .
\]
Then $\Sigma_p$ is diagonal plus rank one,
\begin{equation}\label{eq:cnb-sigma}
  \Sigma_p \;=\; g'_p \diag(u_{p\cdot}) \;+\; g''_p\, u_{p\cdot} u_{p\cdot}^\top ,
\end{equation}
and consequently
\begin{align}
  \frac{\partial^2 f_p}{\partial\beta_{ck}\partial\beta_{c'k'}}
   &= L_{pc}L_{pc'}\Bigl[g'_p\bigl(B^\top\diag(u_{p\cdot})B\bigr)_{kk'}
      + g''_p\, a_{pk}a_{pk'}\Bigr],\\
  \frac{\partial^2 f_p}{\partial\beta_{ck}\partial\gamma_q}
   &= L_{pc}\, a_{pk}\, \sum_{i\in\E_p}\frac{(y_{i\cdot}+\nu_p)\nu_p e_i G_{iq}}{D_i^2},\\
  \frac{\partial^2 f_p}{\partial\gamma_q\partial\gamma_{q'}}
   &= \sum_{i\in\E_p}\frac{(y_{i\cdot}+\nu_p)\nu_p E_p e_i}{D_i^2}\,G_{iq}G_{iq'} .
\end{align}
Only the experiment totals $y_{i\cdot}$ appear; the voxelwise marginals $Y_{pv}$ enter
\eqref{eq:cnb-ll} linearly and so vanish.
\end{proposition}

\begin{proof}
Let $\Gamma_p(E,\gamma) = \sum_{i\in\E_p}(y_{i\cdot}+\nu_p)\log(Ee_i+\nu_p)$, so that
$f_p = \Gamma_p(E_p, \gamma)$ up to terms linear in $\phi$ and terms free of $\phi$. Since
$\partial E_p/\partial S_{pv} = u_{pv}$ and
$\partial^2 E_p/\partial S_{pv}\partial S_{pv'} = u_{pv}\delta_{vv'}$,
\[
  \frac{\partial^2 f_p}{\partial S_{pv}\partial S_{pv'}}
   = \Gamma_p''(E_p)\,u_{pv}u_{pv'} + \Gamma_p'(E_p)\,u_{pv}\delta_{vv'} ,
\]
and $\Gamma_p' = g'_p$, $\Gamma_p'' = g''_p$ by direct differentiation, giving
\eqref{eq:cnb-sigma}. Applying Lemma~\ref{lem:master} and using
$B^\top(u u^\top)B = (B^\top u)(B^\top u)^\top = a_pa_p^\top$ gives the $\beta\beta$ block. For
the cross block, $\partial_q \bigl[e_i/D_i\bigr] = \nu_p e_i G_{iq}/D_i^2$, whence
$\partial_q \Gamma'_p = \sum_i (y_{i\cdot}+\nu_p)\nu_pe_iG_{iq}/D_i^2$ and the
$\beta$-derivative contributes the factor $L_{pc}a_{pk}$. For the $\gamma\gamma$ block,
$\partial_q \Gamma_p = \sum_i (y_{i\cdot}+\nu_p)E_pe_iG_{iq}/D_i$ and
$\partial_{q'}\bigl[E_pe_i/D_i\bigr] = E_pe_i\nu_pG_{iq'}/D_i^2$.
\end{proof}

\begin{remark}
Proposition~\ref{prop:cnb} is the only one of the three whose $\Sigma_p$ is not diagonal, and it
is nonetheless the cheapest to prove: with $\nu_p$ fixed, every special function in
\eqref{eq:cnb-ll} depends on constants alone.
\end{remark}

\section{Equivalence}

Two independent checks establish that the propositions describe the same matrix NiMARE computes.

\subsection{Symbolic}

\texttt{proofs/negative\_binomial.py} and \texttt{proofs/clustered\_negative\_binomial.py}
construct the likelihood exactly as \texttt{nimare/meta/cbmr/distributions.py} writes it,
differentiate it with \texttt{sympy}, subtract the claimed closed form, and simplify. Every
residual is the symbol $0$; nothing numerical is involved, so these are proofs on a symbolic
instance rather than evidence at a sample point.

For the clustered case the whole composite Hessian closes in one pass. For the negative binomial
the composite is intractable for \texttt{simplify}, so the proof is decomposed exactly as the
argument above is: (i) NiMARE's summand equals $\lambda$ of Lemma~\ref{lem:nb-lambda}; (ii) the
partials of $\lambda$ are those of Proposition~\ref{prop:nb}, and $S$ enters as $\eta$ does;
(iii) the derivatives of $\eta,\kappa$ are those of Lemma~\ref{lem:nb-mk-derivs}. Composition is
then Lemma~\ref{lem:master} together with the chain rule, which is a theorem. Two identities the
implementation relies on for precision are checked as well: at integer counts the digamma terms
equal the sums \eqref{eq:nb-count-sums}, and $g$ has the power series of
Remark~\ref{rem:nb-precision}.

\subsection{Numerical}

\texttt{scripts/test\_cbmr\_analytic.py} and \texttt{scripts/test\_cbmr\_designs.py} compare
the implementations --- vectorisation, pattern assembly, loadings scatter --- against
\texttt{CBMRModel.information\_matrix}, i.e. against \texttt{torch.func.hessian} itself.
Table~\ref{tab:designs} reports the design coverage. Two further tests guard the structural
claims: the Poisson Hessian must be unchanged when the foci are doubled
(Proposition~\ref{prop:poisson}), and the other two must move, so a foci-independent
implementation of either cannot pass silently. Near $\theta_p = 0$ autodiff is itself inexact
(Remark~\ref{rem:nb-precision}), so NiMARE's test of the negative binomial there compares against
a $50$-digit \texttt{mpmath} Hessian of the likelihood instead.

\section{Structure of the information matrix}

\subsection{Sparsity from the design}

\begin{proposition}[Column separation]\label{prop:blocks}
Let $\supp(p) = \{c : L_{pc} \neq 0\}$ and let $\sim$ be the transitive closure of the relation
$c \sim c' \iff \exists p:\ c,c' \in \supp(p)$. If $Q = 0$ and $c \not\sim c'$ then
$H_{(ck),(c'k')} = 0$ for all $k,k'$. Consequently $H$ is block diagonal over the connected
components of $\sim$, each block of size $|{\cdot}| \cdot K$.
\end{proposition}

\begin{proof}
Immediate from \eqref{eq:master-bb}: each summand carries the factor $L_{pc}L_{pc'}$, which
vanishes unless both $c$ and $c'$ lie in $\supp(p)$; if they do, then $c \sim c'$.
\end{proof}

Proposition~\ref{prop:blocks} holds for \emph{any} $\Sigma_p$ and hence for all three
distributions --- and indeed for any likelihood of the form \eqref{eq:decomposition}. It is a
statement about the design, not about the observation model. Note the entries are
\emph{structurally} zero, not small: no rounding decision is involved.

\subsection{Which designs separate}

The hypothesis $Q=0$ is not cosmetic, and neither is the connectivity condition. A cell-means
spatial factor \texttt{s(a)} assigns each pattern exactly one column, so $\sim$ is the identity
and there are $C$ blocks; so does a full interaction \texttt{s(a:b)}. Every other construction
the CBMR grammar admits couples the columns:
\begin{itemize}
  \item a sum-to-zero term \texttt{sz(a)} spans all levels in every row, so $\supp(p)$ is the
        whole term and one component absorbs everything;
  \item a spatial covariate \texttt{s(x)} loads alongside whatever factor accompanies it;
  \item any global moderator makes $Q > 0$, and the cross block \eqref{eq:master-bg} is generally
        dense, coupling every spatial column through $\gamma$.
\end{itemize}
Table~\ref{tab:designs} makes this concrete. The practical consequence is that the closed-form
Hessian is general while block-diagonal inversion is not, and the two must be gated separately.

\subsection{Bordered structure}\label{sec:bordered}

When the spatial columns separate but $Q>0$, $H$ is not block diagonal but is \emph{bordered}
block diagonal,
\begin{equation}\label{eq:bordered}
  H = \begin{pmatrix} D & C \\ C^\top & E\end{pmatrix},
  \qquad D = \operatorname{blkdiag}(D_1,\dots,D_M) \in \R^{CK\times CK},\quad C \in \R^{CK\times Q},
\end{equation}
with $Q$ typically between $1$ and $5$.

\begin{proposition}[Bordered inverse]\label{prop:schur}
If $D$ and the Schur complement $\Sigma = E - C^\top D^{-1}C$ are nonsingular, then
\begin{equation}
  H^{-1} =
  \begin{pmatrix}
    D^{-1} + D^{-1}C\,\Sigma^{-1}C^\top D^{-1} & -D^{-1}C\,\Sigma^{-1}\\[2pt]
    -\Sigma^{-1}C^\top D^{-1} & \Sigma^{-1}
  \end{pmatrix},
\end{equation}
computable in $O\bigl(\sum_j |D_j|^3\bigr) + O(CKQ\cdot\max_j|D_j|) + O(Q^3)$ plus the
$O((CK)^2Q)$ cost of forming the outer product, rather than $O((CK+Q)^3)$.
\end{proposition}

\begin{proof}
Standard block inversion; verify by multiplying out $H H^{-1} = I$ blockwise, using
$E - C^\top D^{-1} C = \Sigma$.
\end{proof}

Proposition~\ref{prop:schur} extends the fast \emph{inverse} to designs of the form
``cell-means factor plus global moderators''. It does not extend the fast \emph{condition
number}: the spectrum of \eqref{eq:bordered} is not the union of the spectra of $D$ and
$\Sigma$, and no interlacing bound recovers the value.

\section{Condition number and inversion}

\texttt{CBMRModel.covariance} computes $\cond_2(H) = \sigma_{\max}/\sigma_{\min}$ with
\texttt{np.linalg.cond}, which forms a full singular value decomposition, and then inverts with
\texttt{np.linalg.inv}. At the $21$-group shape these cost $163$~s and $14$~s respectively, so
the diagnostic dominates the quantity it is diagnosing. Three generalisations are available, in
increasing order of applicability.

\begin{proposition}[Block-diagonal spectrum]\label{prop:spectrum}
If $H = \operatorname{blkdiag}(H_1,\dots,H_M)$ then
$\spec(H) = \bigsqcup_j \spec(H_j)$ and the
singular values are likewise the multiset union. Hence
$\cond_2(H) = \max_j\sigma_{\max}(H_j) \big/ \min_j \sigma_{\min}(H_j)$, exactly.
\end{proposition}

\begin{proof}
$H - \lambda I$ is block diagonal with blocks $H_j - \lambda I$, and the determinant of a block
diagonal matrix is the product of the blocks' determinants, so the characteristic polynomial
factorises. The same argument applied to $H^\top H$ gives the singular values.
\end{proof}

\begin{proposition}[Symmetry suffices]\label{prop:eigh}
$H$ is symmetric by construction. For symmetric $H$ with eigenvalues $\lambda_i$, the singular
values are $|\lambda_i|$, so $\cond_2(H) = \max_i|\lambda_i| / \min_i|\lambda_i|$ and a symmetric
eigensolver may replace the SVD. This requires no hypothesis on the design and no positive
definiteness.
\end{proposition}

\begin{proof}
Write $H = V\Lambda V^\top$ with $V$ orthogonal. Then $H^\top H = V\Lambda^2V^\top$, so the
singular values --- the nonnegative square roots of the eigenvalues of $H^\top H$ --- are
$|\lambda_i|$.
\end{proof}

\begin{proposition}[One factorisation for both]\label{prop:chol}
If $H \succ 0$, a single Cholesky factorisation $H = LL^\top$ serves both purposes: LAPACK's
\texttt{pocon} returns an estimate of $\cond_1(H)$ in $O(n^2)$ from $L$, and \texttt{potri}
returns $H^{-1}$ in $O(n^3/3)$ from the same $L$. Moreover $\cond_1$ and $\cond_2$ are equivalent
within a factor of $n$:
\begin{equation}\label{eq:norm-equiv}
  n^{-1}\cond_1(H) \;\le\; \cond_2(H) \;\le\; n\,\cond_1(H).
\end{equation}
\end{proposition}

\begin{proof}
The factorisation and the two LAPACK routines are standard. For \eqref{eq:norm-equiv}, from
$n^{-1/2}\|M\|_1 \le \|M\|_2 \le n^{1/2}\|M\|_1$ applied to $M = H$ and $M = H^{-1}$,
$\cond_2 = \|H\|_2\|H^{-1}\|_2 \le n^{1/2}\|H\|_1 \cdot n^{1/2}\|H^{-1}\|_1 = n\cond_1$, and
symmetrically.
\end{proof}

\begin{remark}[The bound is not the behaviour]\label{rem:pocon}
Proposition~\ref{prop:chol} is the fastest route by a wide margin, and \eqref{eq:norm-equiv}
makes it sound harmless: a factor of $n$ is four of the sixteen decimal digits. Measured, it is
not harmless. On the three structures of Table~\ref{tab:strategies} the \texttt{pocon} estimate
reports $90.09$ against a true $\cond_2$ of $5.02$, $2.42\times10^{5}$ against $1657.7$, and
$1312.4$ against $4.98$ --- discrepancies of $18\times$, $146\times$ and $264\times$. These are
comfortably inside the bound, which is the point: the bound permits behaviour bad enough to
make the estimate useless as a drop-in for a number the user sees. A fit at
$\cond_2 = 10^{14}$ --- poorly conditioned but still invertible --- would be reported at
$\sim\!10^{16}$ and trip a warning saying its standard errors cannot be trusted.

The conclusion is to split the two uses of the factorisation. The Cholesky is adopted for the
\emph{inverse}, where it agrees with \texttt{np.linalg.inv} to $1.8\times10^{-15}$ and is
$2$--$3\times$ faster. It is rejected for the \emph{condition number}, which comes instead from
Proposition~\ref{prop:eigh}: exact, fully general, and still $3$--$4\times$ faster than the SVD
the stock path forms.
\end{remark}

\subsection{The ladder}\label{sec:ladder}

Collecting Propositions~\ref{prop:schur}--\ref{prop:chol} into a dispatch, most specific rung
first:

\begin{enumerate}
  \item \textbf{blockwise} --- the spatial columns separate and $Q = 0$. Condition number by
        Proposition~\ref{prop:spectrum}, inverse blockwise. Reached only by a lone cell-means
        factor or a lone full interaction.
  \item \textbf{bordered} --- the spatial columns separate but $Q > 0$. Inverse by
        Proposition~\ref{prop:schur}; the condition number cannot follow, so it falls back to
        Proposition~\ref{prop:eigh} on the whole matrix.
  \item \textbf{Cholesky} --- no exploitable structure, but $H \succ 0$. Condition number by
        Proposition~\ref{prop:eigh}, inverse by Cholesky.
  \item \textbf{dense} --- $H \not\succ 0$, so the Cholesky declines. Falls through to
        \texttt{np.linalg.inv}, which either succeeds or raises with the diagnostic the stock
        path would have raised.
\end{enumerate}

Every rung reports the \emph{same} condition number, to the last printed digit; none of them is
an estimate. What varies is only how much work is done to obtain it. Structural claims are
checked against the matrix in hand rather than assumed --- a design that ought to separate but
whose information matrix has an entry outside the implied blocks is inverted densely, with a
warning, rather than silently truncated.

Table~\ref{tab:designs} records which rung each design reaches. The fourth is reached by exactly
one design in the suite, \texttt{\textasciitilde\ sz(a) + sz(b) + sz(a:b)}, which is rank
deficient at that sample size ($\cond_2 \approx 4\times10^{20}$); the Cholesky declining is the
guard working.

\section{Empirical results}

\subsection{Cost of the closed forms}

\input{table_scale}

\subsection{Design coverage}

\input{table_designs}

\subsection{Condition number and inverse}

\input{table_strategies}

\section{Scope and limitations}

\begin{enumerate}
  \item The derivations hold with the overdispersion parameters $\theta_p$ held fixed. This
        matches what \texttt{information\_matrix} does and what CBMR reports, but it means the
        reported standard errors do not propagate uncertainty in $\theta$. Extending to the joint
        $(\phi,\theta)$ information would require the $\theta$ blocks, which are not derived here.
  \item NiMARE's \texttt{check\_design} refuses the overdispersed distributions when any pattern
        has a single experiment, which excludes a continuous spatial covariate. That is a
        pre-existing restriction of the moment-matching approximation, not of these derivations;
        the same term with a coarse covariate is accepted and handled.
  \item The negative binomial closed form stays accurate as $\theta_p \to 0$
        (Remark~\ref{rem:nb-precision}), but the likelihood NiMARE fits is still written in
        $r = 1/k_p$. A fit that drives $\theta_p$ toward zero evaluates
        $\log\Gamma(Y+r) - \log\Gamma(r)$ with $r$ of order $10^{12}$, which loses precision in
        the fit and in the autodiff Hessian, though not in the closed form.
  \item The covariance is returned as a dense $(CK+Q)^2$ array because
        \texttt{contrasts.py} indexes it densely. At the $21$-group shape that is $430$~MB, and it
        is the remaining memory floor; exploiting the block structure past the covariance itself
        would require changes inside NiMARE.
\end{enumerate}

\clearpage
\input{generated_proofs}

\end{document}
